% ============================================================
% Приложение А — исходный код имитатора пользовательских заявок
% Файл: appendices/A_code.tex
%
% ИНСТРУКЦИЯ:
%   Вставьте содержимое вашего файла simulator.py между
%   маркерами  <<< НАЧАЛО КОДА >>> и <<< КОНЕЦ КОДА >>>
%   Если файл длинный — разбейте на несколько lstlisting.
% ============================================================

\newpage
\section*{ПРИЛОЖЕНИЕ А. ИСХОДНЫЙ КОД}
\addcontentsline{toc}{section}{ПРИЛОЖЕНИЕ А. ИСХОДНЫЙ КОД}
\renewcommand{\thesection}{А}
 % сбросить счётчик листингов для приложения

В приложении приведён полный исходный код программного имитатора
пользовательских заявок (\texttt{simulator.py}), разработанного на языке
Python~3. Программа выполняет статистический анализ набора учётных
данных SLURM, генерирует синтетические заявки в соответствии с
нормальным и логнормальным законами распределения и автоматически
ставит их в очередь диспетчера посредством утилиты \texttt{sbatch}.

% ------------------------------------------------------------------
% Листинг А.1 — simulator.py (полный текст)
% ------------------------------------------------------------------
\lstinputlisting[
    language=Python,
    caption={Имитатор пользовательских заявок \texttt{simulator.py}},
    label={lst:simulator}
]{simulator.py}